% !PRINTED 456
\weberpage{456}
%<<CONT
scheidung nötig ist, den Körper $\mathfrak{K}(\Delta)$ den \gesperrt{Klassenkörper von $\Omega$}, $\mathfrak{K}(D)$ den \gesperrt{Klassenkörper für die Diskriminante $D$} oder den \gesperrt{Ordnungskörper für den Führer $Q$} nennen.

\gesperrt{Nach den Resultaten von \S\,122 ist, vielleicht mit einer endlichen Zahl von Ausnahmen,}

\begin{quote}
\gesperrt{ein Primideal $\mathfrak{p}$ des Körpers $\Omega$ im Körper $\mathfrak{K}(\Delta)$ dann und nur dann in Primideale ersten Grades zerlegbar, wenn $\mathfrak{p}$ in $\Omega$ vom ersten Grade ist und der Hauptklasse angehört, also gleich einer komplexen Primzahl}
\[
\pi = \frac{x+y\sqrt{\Delta}}{2}
\]
ist.
\end{quote}

Soll $\mathfrak{p}$ in $\mathfrak{K}(D)$ in Primideale ersten Grades $\mathfrak{P}$ zerlegt werden, so kommt noch hinzu, daß $\pi$ der durch $Q$ bestimmten Ordnung angehöre, d.\,h. nach dem Modul $Q$ mit einer rationalen Zahl kongruent sei (\S\,96).

\begin{center}
\rule{2cm}{0.4pt}
\end{center}
